\documentclass[conference]{IEEEtran}
\IEEEoverridecommandlockouts
\usepackage{cite}
\usepackage{amsmath,amssymb,amsfonts}
\usepackage{graphicx}
\usepackage{textcomp}
\usepackage{xcolor}
\usepackage{booktabs}
\usepackage{hyperref}

\def\BibTeX{{\rm B\kern-.05em{\sc i\kern-.025em b}\kern-.08em
    T\kern-.1667em\lower.7ex\hbox{E}\kern-.125emX}}

\begin{document}

\title{Interactive Fluid Simulation Pendant: Embedded Architecture and Physics Core\\
{\footnotesize \textsuperscript{*}UCU Robotics and Embedded Systems -- Defense 1 Project Report (October 2026)}
}

\author{\IEEEauthorblockN{Ivan Bohatyrov}
\IEEEauthorblockA{\textit{Faculty of Applied Sciences} \\
\textit{Ukrainian Catholic University}\\
Lviv, Ukraine \\
bohatyrov.pn@ucu.edu.ua}
\and
\IEEEauthorblockN{Ostap Sukhomlyn}
\IEEEauthorblockA{\textit{Faculty of Applied Sciences} \\
\textit{Ukrainian Catholic University}\\
Lviv, Ukraine \\
sukhomlyn.pn@ucu.edu.ua}
\and
\IEEEauthorblockN{Artem Serhiichuk}
\IEEEauthorblockA{\textit{Faculty of Applied Sciences} \\
\textit{Ukrainian Catholic University}\\
Lviv, Ukraine \\
serhiichuk.pn@ucu.edu.ua}
\and
\IEEEauthorblockN{Andriy Antoniuk}
\IEEEauthorblockA{\textit{Faculty of Applied Sciences} \\
\textit{Ukrainian Catholic University}\\
Lviv, Ukraine \\
a.antoniuk.pn@ucu.edu.ua}
}

\maketitle

\begin{abstract}
This report presents the system architecture, hardware specification, and physics engine for an autonomous interactive fluid simulation pendant. The device features a circular matrix of 240 SMD LEDs ($\varnothing 42$\,mm active PCB) driven by an STM32 Cortex-M microcontroller that displays real-time fluid dynamics responding to physical tilting and acceleration measured by an ultra-low-power 6-axis IMU sensor. For Defense 1, we implemented and validated a lightweight 2D particle physics engine in pure C++17 with distance constraints. The engine requires only 5.95\,KB of static RAM with zero dynamic allocations ($O(1)$ heap) and executes in 0.010\,ms per step ($>100,000$\,FPS on PC, $\approx 1\%$ estimated CPU load on an 84\,MHz ARM Cortex-M4 at 40\,FPS). Furthermore, drawing from direct consultation and hardware handover from the previous year's team (NUCLEO-L432KC), we establish the Charlieplexing display topology (240 discrete SMD 0603 LEDs via 16 GPIOs), reject power-hungry addressable LEDs, calculate a verified power budget ($\approx 6.2$ hours on a 300\,mAh Li-Po battery), specify the complete Bill of Materials, and lay out our development roadmap toward custom PCB manufacturing for Defense 2.
\end{abstract}

\begin{IEEEkeywords}
Embedded systems, real-time simulation, particle physics, Charlieplexing, STM32, IMU, LSM6DS3, wearable device.
\end{IEEEkeywords}

%% =========================================================================
%% SECTION I: INTRODUCTION & PROJECT VISION
%% =========================================================================
\section{Introduction and Requirements}

\subsection{Project Vision}
The goal of this project is to design and build an autonomous, wearable electronic pendant ($\varnothing 45$\,mm, 11.5\,mm thickness) that simulates liquid sloshing inside a sealed circular chamber. An onboard inertial measurement unit (IMU) continuously measures the direction of gravity and user movement. The liquid is rendered onto a custom circular matrix of 240 blue LEDs, creating a responsive physical simulation of fluid behavior.

\subsection{System Requirements}
For a practical wearable prototype, we defined the following engineering requirements:
\begin{itemize}
    \item \textbf{Display \& Responsiveness}: Target refresh rate of 40--60\,FPS to ensure smooth visual motion without perceptible lag.
    \item \textbf{Dimensions}: Outer enclosure diameter $\le 45$\,mm, internal PCB diameter $\le 42$\,mm, total depth $\le 12$\,mm, suitable for wearing as a neck pendant.
    \item \textbf{Battery Life}: Autonomous runtime of $\ge 5$ hours using a compact 300\,mAh Lithium-Polymer (Li-Po) battery rechargeable via USB Type-C.
    \item \textbf{Microcontroller Constraints}: All physics simulation, IMU filtering, and display multiplexing must execute within a low-power microcontroller with $\le 64$\,KB of internal SRAM.
\end{itemize}

\subsection{Past Team Consultation and Practical Insights}
To mitigate hardware and fabrication risks, our team consulted with the previous year's project team (Daryna) and received their hardware development kit, including a NUCLEO-L432KC board and spare components. Key insights incorporated into our engineering plan include:
\begin{enumerate}
    \item \textbf{Reflow Soldering over Manual Iron}: Soldering 240 small SMD LEDs with a conventional soldering iron introduces high risk of solder bridges and thermal damage. Assembly will be carried out using a hot-air rework station (or hot plate) with solder paste and a laser-cut stainless steel stencil.
    \item \textbf{Redundant Backup PCBs}: Fabrication yields for dense two-layer LED matrices can vary; we are ordering five fabricated boards from JLCPCB to maintain at least two fully assembled working units.
    \item \textbf{EDA Workflow}: Utilizing KiCad for parametric circular footprint placement and trace routing.
\end{enumerate}

%% =========================================================================
%% SECTION II: HARDWARE ARCHITECTURE & SYSTEM DESIGN
%% =========================================================================
\section{Hardware Architecture}

\subsection{Microcontroller Strategy (Prototyping vs. Final PCB)}
We adopt a clear two-stage microcontroller strategy:
\begin{itemize}
    \item \textbf{Stage 1 (Laboratory Bring-up \& Prototyping)}: We utilize the \textbf{NUCLEO-L432KC} board (ARM Cortex-M4 @ 80\,MHz, 64\,KB SRAM) obtained from the previous team, alongside a generic STM32F401 BlackPill. This enables immediate testing of basic I2C sensor communications, timer-driven GPIO multiplexing, and algorithm profiling.
    \item \textbf{Stage 2 (Final Custom Circular PCB)}: We selected the standalone \textbf{STM32F401CCU6} (or pin-compatible STM32F411CEU6) in a compact UFQFPN48 package (ARM Cortex-M4 @ 84\,MHz, hardware single-precision FPU, 256\,KB Flash, 64\,KB SRAM). It provides 36 available GPIO lines, abundant for our Charlieplexing matrix, while fitting neatly inside the $\varnothing 42$\,mm board geometry. Firmware is programmed via SWD using an ST-Link V2 dongle.
\end{itemize}

\subsection{Display Architecture: Charlieplexing vs. Addressable LEDs}
Driving 240 LEDs individually without specialized controller ICs is achieved via \textbf{Charlieplexing}, which leverages the tri-state capability of microcontroller I/O pins (\texttt{HIGH}, \texttt{LOW}, and High-Impedance \texttt{Hi-Z}). With $N$ GPIO pins, the maximum addressable LED count is:
\begin{equation}
    L_{\text{max}} = N(N - 1)
\end{equation}
Allocating $N = 16$ GPIO pins yields:
\begin{equation}
    16 \times (16 - 1) = 240 \text{ LEDs}
\end{equation}

\textbf{Rejection of Addressable LEDs}: We evaluated smart addressable RGB LEDs (such as WS2812B-2020), but rejected them due to severe power penalties: each integrated driver chip consumes $\approx 1$\,mA quiescent current even when completely dark, totaling $\approx 240$\,mA idle draw for 240 LEDs. This would exhaust a 300\,mAh battery in under 50 minutes without even illuminating the matrix. Furthermore, addressable LEDs cannot use Charlieplexing pin reduction. 

Instead, discrete \textbf{SMD 0603 Blue LEDs} were selected. Because Charlieplexing rapidly scans through rows/pairs, only a small fraction of LEDs conduct current at any microsecond, keeping average matrix consumption at a modest $\approx 25$\,mA. Dedicated SMD current-limiting resistors are placed on each of the 16 driving lines to protect MCU GPIO pins.

\subsection{Sensor Subsystem}
Orientation and motion are captured via an \textbf{LSM6DS3} 6-axis ultra-low-power IMU over I2C at 400\,kHz. With an active current draw of only 0.9\,mA, it reduces power consumption significantly compared to older sensors. The accelerometer samples the 2D gravity vector at 104\,Hz, driving the physics simulation frame.

\subsection{Power Budget and Battery Life}
Table~\ref{tab:power} summarizes the continuous current draw at 3.3\,V.
\begin{table}[htbp]
\caption{Estimated System Power Consumption}
\label{tab:power}
\centering
\resizebox{\columnwidth}{!}{%
\begin{tabular}{lcr}
\toprule
\textbf{Subsystem} & \textbf{Operating Mode} & \textbf{Current (mA)} \\
\midrule
STM32F401 MCU             & 84\,MHz active, FPU on         & 15.00\,mA \\
LED Matrix (0603 Blue)    & Charlieplexed dynamic scanning & 25.00\,mA \\
LSM6DS3 IMU               & Accelerometer active (104\,Hz) &  0.90\,mA \\
LDO Regulator (3.3V)      & XC6206 quiescent loss          &  0.05\,mA \\
TP4056 Protection Circuit & Standby discharge leakage      &  0.005\,mA \\
\midrule
\textbf{Total Draw}       &                                & \textbf{40.95\,mA} \\
\bottomrule
\end{tabular}%
}
\end{table}

With an average current of $\approx 41$\,mA, a single 3.7\,V, 300\,mAh Li-Po battery delivers:
\begin{equation}
    t_{\text{run}} \approx \frac{300\,\text{mAh} \times 0.85}{41\,\text{mA}} \approx 6.22 \text{ hours}
\end{equation}
comfortably exceeding the 5-hour target.

%% =========================================================================
%% SECTION III: FLUID SIMULATION PHYSICS CORE
%% =========================================================================
\section{Fluid Simulation Engine}

\subsection{Algorithm Selection Rationale}
Simulating liquid on an embedded microcontroller requires balancing visual realism with strict memory and computation limits. We analyzed three common approaches:
\begin{itemize}
    \item \textbf{Grid-Based Navier-Stokes}: Requires large 2D arrays for velocity and pressure, along with iterative linear solvers. This exceeds our 64\,KB SRAM limit and is computationally prohibitive.
    \item \textbf{Cellular Automata (Falling Sand)}: Very low CPU usage, but particles move only by discrete falling rules without true velocity or momentum. It looks like dry sand rather than flowing water.
    \item \textbf{Particle Simulation with Distance Constraints (Our Choice)}: Liquid is modeled as a small set of discrete 2D particles with geometric non-penetration constraints. It provides continuous momentum, splashes, and wave sloshing while requiring under 6\,KB of RAM and running in microseconds.
\end{itemize}

\subsection{Simulation Equations}
The fluid is represented by $N$ particles inside a circle of normalized radius $R = 1.0$. Each particle has position $\mathbf{x}_i \in \mathbb{R}^2$ and velocity $\mathbf{v}_i \in \mathbb{R}^2$. At each time step $\Delta t$, the engine executes four simple steps:

\subsubsection{1. Motion Prediction}
Particles are updated under velocity and external acceleration $\mathbf{g}$ measured by the IMU:
\begin{align}
    \mathbf{v}_i^* &= \alpha \cdot \mathbf{v}_i + \mathbf{g} \Delta t \\
    \mathbf{x}_i^* &= \mathbf{x}_i + \mathbf{v}_i^* \Delta t
\end{align}
where $\alpha = 0.988$ is a damping factor modeling slight air resistance.

\subsubsection{2. Incompressibility (Particle Spacing)}
To keep liquid volume constant, particles must not overlap. If the distance between two particles $d_{ij} = |\mathbf{x}_i^* - \mathbf{x}_j^*|$ is less than their diameter $2r$, they are pushed apart along their contact normal $\mathbf{n}_{ij}$:
\begin{equation}
    \Delta \mathbf{x}_i = \frac{1}{2} k \, (2r - d_{ij}) \, \mathbf{n}_{ij}, \quad \Delta \mathbf{x}_j = -\Delta \mathbf{x}_i
\end{equation}
where $k = 0.92$ is the collision stiffness.

\subsubsection{3. Circular Boundary Constraint}
Particles are confined inside the circular chamber:
\begin{equation}
    \mathbf{x}_i^* = R_{\text{max}} \frac{\mathbf{x}_i^*}{|\mathbf{x}_i^*|}, \quad \text{if } |\mathbf{x}_i^*| > R_{\text{max}}
\end{equation}
where $R_{\text{max}} = 0.97$.

\subsubsection{4. Velocity Update from Displacement}
Velocity is updated directly from the displacement:
\begin{equation}
    \mathbf{v}_i = \frac{\mathbf{x}_i^* - \mathbf{x}_i}{\Delta t}
\end{equation}
When particles hit the bottom of the container, the collision push forces them laterally. Because velocity is derived from this lateral movement, the impact naturally produces horizontal splash waves without artificial forces.

\subsubsection{5. Neighbor Viscosity Filter (XSPH)}
To ensure particles flow together as a cohesive liquid instead of loose dust, a velocity smoothing filter is applied:
\begin{equation}
    \Delta \mathbf{v}_i = c \sum_{j} (\mathbf{v}_j - \mathbf{v}_i) \left(1 - \frac{d_{ij}}{h}\right)
\end{equation}
where $c = 0.015$ is the viscosity coefficient and $h$ is the neighbor radius.

\subsection{Performance and Benchmark Results}
The engine was implemented in pure embedded C++17 with zero heap usage ($O(1)$ memory). Table~\ref{tab:benchmarks} presents benchmark results measured across different fluid fill levels.

\begin{table}[htbp]
\caption{Simulation Performance across Fluid Fill Levels}
\label{tab:benchmarks}
\centering
\resizebox{\columnwidth}{!}{%
\begin{tabular}{lccccc}
\toprule
\textbf{Fill Level} & \textbf{Particles ($N$)} & \textbf{Active LEDs} & \textbf{Step Time} & \textbf{PC FPS} & \textbf{MCU Load$^*$} \\
\midrule
20\%  & 8   & 36  & 0.002\,ms & 415,000 & $<1\%$ \\
40\%  & 18  & 69  & 0.006\,ms & 170,000 & $<1\%$ \\
\textbf{60\% (Nominal)} & \textbf{28} & \textbf{95} & \textbf{0.010\,ms} & \textbf{103,000} & \textbf{$\approx 1\%$} \\
80\%  & 40  & 125 & 0.016\,ms & 64,000  & $\approx 2\%$ \\
100\% (Full) & 52  & 150 & 0.023\,ms & 43,000  & $\approx 3\%$ \\
\bottomrule
\multicolumn{6}{l}{$^*$Estimated CPU load on ARM Cortex-M4 @ 84\,MHz at 40\,FPS target.}
\end{tabular}%
}
\end{table}

\begin{itemize}
    \item \textbf{SRAM Footprint}: Total simulation state occupies \textbf{5.95\,KB} (9.3\% of the STM32F401 64\,KB RAM).
    \item \textbf{Stability}: Over 10,000 steps subjected to 5G acceleration shocks and rapid rotation, zero particles escaped the boundary and no arithmetic NaN errors occurred.
\end{itemize}

%% =========================================================================
%% SECTION IV: BILL OF MATERIALS & PROCUREMENT
%% =========================================================================
\section{Bill of Materials and Procurement}

Table~\ref{tab:bom} lists the key components required for prototyping and the final pendant device.

\begin{table}[htbp]
\caption{Bill of Materials (BOM)}
\label{tab:bom}
\centering
\resizebox{\columnwidth}{!}{%
\begin{tabular}{llcc}
\toprule
\textbf{Component} & \textbf{Model / Spec} & \textbf{Qty} & \textbf{Unit Cost} \\
\midrule
Lab Dev Board          & NUCLEO-L432KC / BlackPill          & 1     & \$0.00 (Supplied) \\
Standalone MCU         & STM32F401CCU6 (UFQFPN48)           & 1     & \$3.40 \\
IMU Sensor             & LSM6DS3TR-C (6-axis, ultra-LP)     & 1     & \$2.20 \\
Matrix LEDs            & SMD 0603 Blue High-Brightness      & 500   & \$3.00 \\
Current-Limit Resistors& SMD 0603 (16 lines + spares)       & 20    & \$0.50 \\
Decoupling Capacitors  & SMD 0603 (100nF, 1uF, 10uF)        & 1 kit & \$0.75 \\
USB Type-C Connector   & 6-pin SMD Receptacle (Power)       & 2     & \$0.60 \\
Battery Charger IC     & TP4056 with Protection (DW01A)     & 1     & \$0.40 \\
LDO Voltage Regulator  & XC6206P332MR (3.3V, 250mA)         & 2     & \$0.40 \\
Li-Po Battery          & 3.7\,V 300\,mAh (Model 502030)     & 1     & \$2.70 \\
Custom Circular PCB    & 2-layer $\varnothing 42$\,mm (JLCPCB) & 5  & \$8.50 \\
SMD Solder Stencil     & Stainless steel frameless stencil  & 1     & \$3.00 \\
External Programmer    & ST-Link V2 USB Dongle (SWD)        & 1     & \$2.70 \\
\midrule
\textbf{Total Prototype Cost} & & & \textbf{$\approx \$28.15$} \\
\bottomrule
\end{tabular}%
}
\end{table}

%% =========================================================================
%% SECTION V: ENCLOSURE DESIGN
%% =========================================================================
\section{Enclosure and Mechanical Design}
The mechanical enclosure is designed in CAD (Fusion 360) as a two-piece circular puck ($\varnothing 45$\,mm outer diameter, 11.5\,mm depth) accommodating an internal $\varnothing 42$\,mm custom PCB. It consists of:
\begin{itemize}
    \item A 3D-printed bottom housing containing the Li-Po battery pocket (30$\times$20$\times$5\,mm), PCB mounting snap-fits, and a cutout for the 6-pin USB Type-C charging port.
    \item A semi-translucent frosted acrylic diffuser front cover to blend individual LED points into a smooth fluid surface.
    \item An integrated lanyard loop on top for wearing the pendant on a cord or chain.
\end{itemize}
Prototypes will be fabricated using SLA resin 3D printing for fine detail and smooth surface finish.

%% =========================================================================
%% SECTION VI: DEVELOPMENT ROADMAP
%% =========================================================================
\section{Development Roadmap to Defense 2}
The milestones toward the final working device are scheduled as follows:
\begin{itemize}
    \item \textbf{Milestone 1 (Defense 1, Mid-October)}: C++ simulation core verified; Charlieplexing logic validated with NUCLEO-L432KC and breadboard LEDs; BOM and power budget established.
    \item \textbf{Milestone 2 (End of October)}: Complete KiCad schematic and circular PCB layout with 240 Charlieplexed LEDs ($\varnothing 42$\,mm); order 5x PCB batch with stencil from JLCPCB.
    \item \textbf{Milestone 3 (November)}: Solder and assemble two circular PCBs using hot-air gun and solder paste; flash firmware via SWD; integrate LSM6DS3 IMU readings with the simulation loop over I2C.
    \item \textbf{Milestone 4 (Defense 2, Early December)}: Final assembly inside the 3D-printed enclosure; power management tuning; battery endurance testing.
\end{itemize}

%% =========================================================================
%% REFERENCES
%% =========================================================================
\begin{thebibliography}{00}
\bibitem{macklin2013} M. Macklin and M. M{\"u}ller, ``Position based fluids,'' \textit{ACM Transactions on Graphics (TOG)}, vol. 32, no. 4, pp. 1--12, 2013.
\bibitem{charlieplexing} M. Lancaster, ``Charlieplexing: Reduced pin-count LED displays,'' \textit{Microchip Technology Inc. Application Note AN188}, 2001.
\bibitem{muller2003} M. M{\"u}ller, D. Charypar, and M. Gross, ``Particle-based fluid simulation for interactive applications,'' in \textit{Proc. 2003 ACM SIGGRAPH/Eurographics Symposium on Computer Animation}, 2003, pp. 154--164.
\end{thebibliography}

\end{document}
